% --- Archon physics formalization source begin ---
% archon:physics
% archon:covers PhyXMiniProblems/problem_phyx_mini_0308.lean
% archon:source-report reports/phyx_mini/problem_phyx_mini_0308.source.json

\chapter{Physics problem phyx\_mini\_0308}
\label{ch:PhyXMiniProblems_problem_phyx_mini_0308}

\paragraph{Problem source.}
\#\# Physical scenario

One clue used by your brain to determine the direction of a source of sound is the time delay \$\textbackslash{}Delta t\$ between the arrival of the sound at the ear closer to the source and the arrival at the farther ear. Assume that the source is distant so that a wavefront from it is approximately planar when it reaches you, and let \$D\$ represent the separation between your ears.

\#\# Question

Based on the time-delay clue, your brain interprets the submerged sound to arrive at an angle \$\textbackslash{}theta\$ from the forward direction. Evaluate \$\textbackslash{}theta\$ for fresh water at \$20\textasciicircum{}\textbackslash{}circ\$C.

\#\# Answer choices

- A: A: \$4\textasciicircum{}\textbackslash{}circ\$

- B: B: \$7\textasciicircum{}\textbackslash{}circ\$

- C: C: \$10\textasciicircum{}\textbackslash{}circ\$

- D: D: \$13\textasciicircum{}\textbackslash{}circ\$

\#\# Figure caption (auxiliary; use the image as primary evidence)

The image depicts a wavefront diagram with geometric components:

- Two dark points labeled \textbackslash{}(L\textbackslash{}) and \textbackslash{}(R\textbackslash{}) are connected by lines, forming a right triangle with a hypotenuse labeled \textbackslash{}(d\textbackslash{}).
- The angle \textbackslash{}( \textbackslash{}theta \textbackslash{}) is marked at the points \textbackslash{}(L\textbackslash{}) and \textbackslash{}(R\textbackslash{}).
- The line labeled \textbackslash{}(D\textbackslash{}) is perpendicular to the wavefront lines, intersecting the hypotenuse at right angles.
- Two parallel lines with wavefront arrows point towards the hypotenuse at angles marked by \textbackslash{}( \textbackslash{}theta \textbackslash{}).
- The word "Wavefronts" is placed near the parallel lines.

This diagram likely illustrates concepts related to wave propagation, such as reflection or refraction in physics.

\#\# Dataset metadata

Wave Properties; Physical Model Grounding Reasoning, Spatial Relation Reasoning

\paragraph{Recorded answer/context.}
D: D: \$13\textasciicircum{}\textbackslash{}circ\$

\paragraph{Figure/image path.}
/root/proposal\_for\_physic/science-mango/phyx\_mini\_run/phyx\_data/test\_image/308.png

\paragraph{Formalization target.}
create a compiling Lean file with sorry bodies at `PhyXMiniProblems/problem\_phyx\_mini\_0308.lean`. The Lean declarations must preserve the physical quantities, dimensions or dimensional roles, figure labels, governing-law hypotheses, and final relation expressed by this problem.
Use Mathlib/Physlib names found through LeanExplore where available. If a physics API is missing, introduce faithful local abstractions rather than scalar placeholder aliases.

\begin{theorem}[Physics formalization target]
\label{thm:physics:phyx_mini_0308:target}
\lean{PhyXMiniProblems.ProblemPhyXMini0308.submergedSound_interpretedAngle_is_thirteen_degrees}
\uses{def:physics:phyx-mini-0308:phyxminiproblems-problemphyxmini0308-radianstodegrees, def:physics:phyx-mini-0308:phyxminiproblems-problemphyxmini0308-soundlocalizationsetup, def:physics:phyx-mini-0308:phyxminiproblems-problemphyxmini0308-matchessourcefigure, def:physics:phyx-mini-0308:phyxminiproblems-problemphyxmini0308-matchesproblemstatement, def:physics:phyx-mini-0308:phyxminiproblems-problemphyxmini0308-matchessoundspeeddata, def:physics:phyx-mini-0308:phyxminiproblems-problemphyxmini0308-matchesintendedsideonarrival, def:physics:phyx-mini-0308:phyxminiproblems-problemphyxmini0308-hasphysicalparameters, def:physics:phyx-mini-0308:phyxminiproblems-problemphyxmini0308-satisfiessoundlocalizationphysics, def:physics:phyx-mini-0308:phyxminiproblems-problemphyxmini0308-isclosestanglechoice}
The assigned autoformalize agent should translate this physics problem into Lean declarations in the covered file, with theorem and lemma proof bodies written as `by sorry`.
\end{theorem}
\begin{proof}
This is an autoformalization task, not a proof task. Produce statements that can later be proved by the physics prover without weakening the physical model.
\end{proof}

% --- Archon declaration topology begin ---
\section{Declaration topology}
\begin{definition}[Acoustic Length]
  \label{def:physics:phyx-mini-0308:phyxminiproblems-problemphyxmini0308-acousticlength}
  \lean{PhyXMiniProblems.ProblemPhyXMini0308.AcousticLength}
  A nonnegative physical one-dimensional distance
\end{definition}
\begin{proof}
  This is the declared model definition of Acoustic Length.
\end{proof}
\begin{definition}[Acoustic Duration]
  \label{def:physics:phyx-mini-0308:phyxminiproblems-problemphyxmini0308-acousticduration}
  \lean{PhyXMiniProblems.ProblemPhyXMini0308.AcousticDuration}
  A nonnegative physical arrival-time delay
\end{definition}
\begin{proof}
  This is the declared model definition of Acoustic Duration.
\end{proof}
\begin{definition}[Acoustic Speed]
  \label{def:physics:phyx-mini-0308:phyxminiproblems-problemphyxmini0308-acousticspeed}
  \lean{PhyXMiniProblems.ProblemPhyXMini0308.AcousticSpeed}
  A nonnegative physical propagation speed
\end{definition}
\begin{proof}
  This is the declared model definition of Acoustic Speed.
\end{proof}
\begin{definition}[length In Meters]
  \label{def:physics:phyx-mini-0308:phyxminiproblems-problemphyxmini0308-lengthinmeters}
  \lean{PhyXMiniProblems.ProblemPhyXMini0308.lengthInMeters}
  \uses{def:physics:phyx-mini-0308:phyxminiproblems-problemphyxmini0308-acousticlength}
  Metre readout of a physical acoustic length
\end{definition}
\begin{proof}
  This is the declared model definition of length In Meters.
\end{proof}
\begin{definition}[duration In Seconds]
  \label{def:physics:phyx-mini-0308:phyxminiproblems-problemphyxmini0308-durationinseconds}
  \lean{PhyXMiniProblems.ProblemPhyXMini0308.durationInSeconds}
  \uses{def:physics:phyx-mini-0308:phyxminiproblems-problemphyxmini0308-acousticduration}
  Second readout of a physical acoustic duration
\end{definition}
\begin{proof}
  This is the declared model definition of duration In Seconds.
\end{proof}
\begin{definition}[speed In Meters Per Second]
  \label{def:physics:phyx-mini-0308:phyxminiproblems-problemphyxmini0308-speedinmeterspersecond}
  \lean{PhyXMiniProblems.ProblemPhyXMini0308.speedInMetersPerSecond}
  \uses{def:physics:phyx-mini-0308:phyxminiproblems-problemphyxmini0308-acousticspeed}
  Metre-per-second readout of a physical acoustic speed
\end{definition}
\begin{proof}
  This is the declared model definition of speed In Meters Per Second.
\end{proof}
\begin{definition}[radians To Degrees]
  \label{def:physics:phyx-mini-0308:phyxminiproblems-problemphyxmini0308-radianstodegrees}
  \lean{PhyXMiniProblems.ProblemPhyXMini0308.radiansToDegrees}
  Convert a radian readout to a degree readout
\end{definition}
\begin{proof}
  This is the declared model definition of radians To Degrees.
\end{proof}
\begin{definition}[Ear Point Label]
  \label{def:physics:phyx-mini-0308:phyxminiproblems-problemphyxmini0308-earpointlabel}
  \lean{PhyXMiniProblems.ProblemPhyXMini0308.EarPointLabel}
  Labels printed beside the two ear points
\end{definition}
\begin{proof}
  This is the declared model definition of Ear Point Label.
\end{proof}
\begin{definition}[Distance Label]
  \label{def:physics:phyx-mini-0308:phyxminiproblems-problemphyxmini0308-distancelabel}
  \lean{PhyXMiniProblems.ProblemPhyXMini0308.DistanceLabel}
  Labels printed on the baseline and perpendicular offset segments
\end{definition}
\begin{proof}
  This is the declared model definition of Distance Label.
\end{proof}
\begin{definition}[Bearing Label]
  \label{def:physics:phyx-mini-0308:phyxminiproblems-problemphyxmini0308-bearinglabel}
  \lean{PhyXMiniProblems.ProblemPhyXMini0308.BearingLabel}
  The angular symbol printed twice in the image
\end{definition}
\begin{proof}
  This is the declared model definition of Bearing Label.
\end{proof}
\begin{definition}[Bearing Meaning]
  \label{def:physics:phyx-mini-0308:phyxminiproblems-problemphyxmini0308-bearingmeaning}
  \lean{PhyXMiniProblems.ProblemPhyXMini0308.BearingMeaning}
  Geometric meanings of the two displayed copies of theta
\end{definition}
\begin{proof}
  This is the declared model definition of Bearing Meaning.
\end{proof}
\begin{definition}[Wavefront Arrangement]
  \label{def:physics:phyx-mini-0308:phyxminiproblems-problemphyxmini0308-wavefrontarrangement}
  \lean{PhyXMiniProblems.ProblemPhyXMini0308.WavefrontArrangement}
  Relationship between the two drawn wavefront lines
\end{definition}
\begin{proof}
  This is the declared model definition of Wavefront Arrangement.
\end{proof}
\begin{definition}[Intersection Kind]
  \label{def:physics:phyx-mini-0308:phyxminiproblems-problemphyxmini0308-intersectionkind}
  \lean{PhyXMiniProblems.ProblemPhyXMini0308.IntersectionKind}
  The marked intersection of d with a wavefront
\end{definition}
\begin{proof}
  This is the declared model definition of Intersection Kind.
\end{proof}
\begin{definition}[Propagation Arrow Sense]
  \label{def:physics:phyx-mini-0308:phyxminiproblems-problemphyxmini0308-propagationarrowsense}
  \lean{PhyXMiniProblems.ProblemPhyXMini0308.PropagationArrowSense}
  Sense of the arrows normal to the wavefronts
\end{definition}
\begin{proof}
  This is the declared model definition of Propagation Arrow Sense.
\end{proof}
\begin{definition}[Source Figure]
  \label{def:physics:phyx-mini-0308:phyxminiproblems-problemphyxmini0308-sourcefigure}
  \lean{PhyXMiniProblems.ProblemPhyXMini0308.SourceFigure}
  \uses{def:physics:phyx-mini-0308:phyxminiproblems-problemphyxmini0308-earpointlabel, def:physics:phyx-mini-0308:phyxminiproblems-problemphyxmini0308-distancelabel, def:physics:phyx-mini-0308:phyxminiproblems-problemphyxmini0308-bearinglabel, def:physics:phyx-mini-0308:phyxminiproblems-problemphyxmini0308-bearingmeaning, def:physics:phyx-mini-0308:phyxminiproblems-problemphyxmini0308-wavefrontarrangement, def:physics:phyx-mini-0308:phyxminiproblems-problemphyxmini0308-intersectionkind, def:physics:phyx-mini-0308:phyxminiproblems-problemphyxmini0308-propagationarrowsense}
  A structured transcription of the primary raster. It records only labels and qualitative incidence data; the image supplies no numerical length or angle
\end{definition}
\begin{proof}
  This is the declared model definition of Source Figure.
\end{proof}
\begin{definition}[Source Distance Regime]
  \label{def:physics:phyx-mini-0308:phyxminiproblems-problemphyxmini0308-sourcedistanceregime}
  \lean{PhyXMiniProblems.ProblemPhyXMini0308.SourceDistanceRegime}
  The distance regime stated in the prose
\end{definition}
\begin{proof}
  This is the declared model definition of Source Distance Regime.
\end{proof}
\begin{definition}[Sound Medium]
  \label{def:physics:phyx-mini-0308:phyxminiproblems-problemphyxmini0308-soundmedium}
  \lean{PhyXMiniProblems.ProblemPhyXMini0308.SoundMedium}
  Medium whose speed is used for actual or interpreted propagation
\end{definition}
\begin{proof}
  This is the declared model definition of Sound Medium.
\end{proof}
\begin{definition}[Sound Localization Setup]
  \label{def:physics:phyx-mini-0308:phyxminiproblems-problemphyxmini0308-soundlocalizationsetup}
  \lean{PhyXMiniProblems.ProblemPhyXMini0308.SoundLocalizationSetup}
  \uses{def:physics:phyx-mini-0308:phyxminiproblems-problemphyxmini0308-acousticlength, def:physics:phyx-mini-0308:phyxminiproblems-problemphyxmini0308-acousticduration, def:physics:phyx-mini-0308:phyxminiproblems-problemphyxmini0308-acousticspeed, def:physics:phyx-mini-0308:phyxminiproblems-problemphyxmini0308-earpointlabel, def:physics:phyx-mini-0308:phyxminiproblems-problemphyxmini0308-sourcefigure, def:physics:phyx-mini-0308:phyxminiproblems-problemphyxmini0308-sourcedistanceregime, def:physics:phyx-mini-0308:phyxminiproblems-problemphyxmini0308-soundmedium}
  Physical data for one localization event. The actual submerged path difference and the air-calibrated inferred path difference are independent dimensionful quantities. The interpreted bearing is also an independent observation; no requested numerical answer is built into this record. The final four real fields quantify the far-field approximation rather than silently identifying an approximately planar wave with an exact plane wave
\end{definition}
\begin{proof}
  This is the declared model definition of Sound Localization Setup.
\end{proof}
\begin{definition}[Matches Source Figure]
  \label{def:physics:phyx-mini-0308:phyxminiproblems-problemphyxmini0308-matchessourcefigure}
  \lean{PhyXMiniProblems.ProblemPhyXMini0308.MatchesSourceFigure}
  \uses{def:physics:phyx-mini-0308:phyxminiproblems-problemphyxmini0308-soundlocalizationsetup}
  Label and incidence information read from 308.png. The primary raster, not the auxiliary caption, shows D along the ear baseline and lowercase d perpendicular to a wavefront. For the arrow orientation, R is reached before L
\end{definition}
\begin{proof}
  This is the declared model definition of Matches Source Figure.
\end{proof}
\begin{definition}[Matches Problem Statement]
  \label{def:physics:phyx-mini-0308:phyxminiproblems-problemphyxmini0308-matchesproblemstatement}
  \lean{PhyXMiniProblems.ProblemPhyXMini0308.MatchesProblemStatement}
  \uses{def:physics:phyx-mini-0308:phyxminiproblems-problemphyxmini0308-soundlocalizationsetup}
  Distant-source and fresh-water data stated in the problem
\end{definition}
\begin{proof}
  This is the declared model definition of Matches Problem Statement.
\end{proof}
\begin{definition}[Matches Sound Speed Data]
  \label{def:physics:phyx-mini-0308:phyxminiproblems-problemphyxmini0308-matchessoundspeeddata}
  \lean{PhyXMiniProblems.ProblemPhyXMini0308.MatchesSoundSpeedData}
  \uses{def:physics:phyx-mini-0308:phyxminiproblems-problemphyxmini0308-speedinmeterspersecond, def:physics:phyx-mini-0308:phyxminiproblems-problemphyxmini0308-soundlocalizationsetup}
  Standard textbook sound-speed readouts at 20 °C: 344 m/s for the listener's ordinary air calibration and 1482 m/s for fresh water
\end{definition}
\begin{proof}
  This is the declared model definition of Matches Sound Speed Data.
\end{proof}
\begin{definition}[Matches Intended Side On Arrival]
  \label{def:physics:phyx-mini-0308:phyxminiproblems-problemphyxmini0308-matchesintendedsideonarrival}
  \lean{PhyXMiniProblems.ProblemPhyXMini0308.MatchesIntendedSideOnArrival}
  \uses{def:physics:phyx-mini-0308:phyxminiproblems-problemphyxmini0308-soundlocalizationsetup}
  The recorded 13° choice requires the submerged source to be directly to one side, giving the maximum interaural delay. The extracted prose does not state this datum, so it is exposed as a separate intended-scenario premise rather than hidden in a law or in the answer
\end{definition}
\begin{proof}
  This is the declared model definition of Matches Intended Side On Arrival.
\end{proof}
\begin{definition}[Has Physical Parameters]
  \label{def:physics:phyx-mini-0308:phyxminiproblems-problemphyxmini0308-hasphysicalparameters}
  \lean{PhyXMiniProblems.ProblemPhyXMini0308.HasPhysicalParameters}
  \uses{def:physics:phyx-mini-0308:phyxminiproblems-problemphyxmini0308-lengthinmeters, def:physics:phyx-mini-0308:phyxminiproblems-problemphyxmini0308-durationinseconds, def:physics:phyx-mini-0308:phyxminiproblems-problemphyxmini0308-speedinmeterspersecond, def:physics:phyx-mini-0308:phyxminiproblems-problemphyxmini0308-soundlocalizationsetup}
  Positivity and principal-angle conditions for the physical setup
\end{definition}
\begin{proof}
  This is the declared model definition of Has Physical Parameters.
\end{proof}
\begin{definition}[Satisfies Bounded Far Field Approximation]
  \label{def:physics:phyx-mini-0308:phyxminiproblems-problemphyxmini0308-satisfiesboundedfarfieldapproximation}
  \lean{PhyXMiniProblems.ProblemPhyXMini0308.SatisfiesBoundedFarFieldApproximation}
  \uses{def:physics:phyx-mini-0308:phyxminiproblems-problemphyxmini0308-lengthinmeters, def:physics:phyx-mini-0308:phyxminiproblems-problemphyxmini0308-soundlocalizationsetup}
  Quantitative meaning of “approximately planar.” The actual finite-distance path difference equals the plane-wave first term plus a signed residual. Its absolute error is bounded, and both the baseline-to-distance ratio and the relative error are controlled by the stored tolerance. Thus no far-field approximation is globalized into an exact equality
\end{definition}
\begin{proof}
  This is the declared model definition of Satisfies Bounded Far Field Approximation.
\end{proof}
\begin{definition}[Satisfies Exact Side On Source Geometry]
  \label{def:physics:phyx-mini-0308:phyxminiproblems-problemphyxmini0308-satisfiesexactsideonsourcegeometry}
  \lean{PhyXMiniProblems.ProblemPhyXMini0308.SatisfiesExactSideOnSourceGeometry}
  \uses{def:physics:phyx-mini-0308:phyxminiproblems-problemphyxmini0308-soundlocalizationsetup}
  Exact finite-source geometry for a source on the extension of the ear baseline. At a side-on bearing, the farther source-to-ear distance exceeds the nearer one by exactly the ear separation; this fact does not use the far-field approximation
\end{definition}
\begin{proof}
  This is the declared model definition of Satisfies Exact Side On Source Geometry.
\end{proof}
\begin{definition}[Satisfies Water Delay Law]
  \label{def:physics:phyx-mini-0308:phyxminiproblems-problemphyxmini0308-satisfieswaterdelaylaw}
  \lean{PhyXMiniProblems.ProblemPhyXMini0308.SatisfiesWaterDelayLaw}
  \uses{def:physics:phyx-mini-0308:phyxminiproblems-problemphyxmini0308-lengthinmeters, def:physics:phyx-mini-0308:phyxminiproblems-problemphyxmini0308-durationinseconds, def:physics:phyx-mini-0308:phyxminiproblems-problemphyxmini0308-speedinmeterspersecond, def:physics:phyx-mini-0308:phyxminiproblems-problemphyxmini0308-soundlocalizationsetup}
  Constant-speed water propagation converts actual path difference to delay
\end{definition}
\begin{proof}
  This is the declared model definition of Satisfies Water Delay Law.
\end{proof}
\begin{definition}[Satisfies Air Calibrated Delay Distance Law]
  \label{def:physics:phyx-mini-0308:phyxminiproblems-problemphyxmini0308-satisfiesaircalibrateddelaydistancelaw}
  \lean{PhyXMiniProblems.ProblemPhyXMini0308.SatisfiesAirCalibratedDelayDistanceLaw}
  \uses{def:physics:phyx-mini-0308:phyxminiproblems-problemphyxmini0308-lengthinmeters, def:physics:phyx-mini-0308:phyxminiproblems-problemphyxmini0308-durationinseconds, def:physics:phyx-mini-0308:phyxminiproblems-problemphyxmini0308-speedinmeterspersecond, def:physics:phyx-mini-0308:phyxminiproblems-problemphyxmini0308-soundlocalizationsetup}
  The auditory system converts the measured delay to a distance using its air sound-speed calibration
\end{definition}
\begin{proof}
  This is the declared model definition of Satisfies Air Calibrated Delay Distance Law.
\end{proof}
\begin{definition}[Satisfies Interpreted Bearing Geometry]
  \label{def:physics:phyx-mini-0308:phyxminiproblems-problemphyxmini0308-satisfiesinterpretedbearinggeometry}
  \lean{PhyXMiniProblems.ProblemPhyXMini0308.SatisfiesInterpretedBearingGeometry}
  \uses{def:physics:phyx-mini-0308:phyxminiproblems-problemphyxmini0308-lengthinmeters, def:physics:phyx-mini-0308:phyxminiproblems-problemphyxmini0308-soundlocalizationsetup}
  The exact right triangle in the interpretation schematic relates the inferred distance d, ear separation D, and the interpreted principal bearing
\end{definition}
\begin{proof}
  This is the declared model definition of Satisfies Interpreted Bearing Geometry.
\end{proof}
\begin{definition}[Satisfies Sound Localization Physics]
  \label{def:physics:phyx-mini-0308:phyxminiproblems-problemphyxmini0308-satisfiessoundlocalizationphysics}
  \lean{PhyXMiniProblems.ProblemPhyXMini0308.SatisfiesSoundLocalizationPhysics}
  \uses{def:physics:phyx-mini-0308:phyxminiproblems-problemphyxmini0308-soundlocalizationsetup, def:physics:phyx-mini-0308:phyxminiproblems-problemphyxmini0308-satisfiesboundedfarfieldapproximation, def:physics:phyx-mini-0308:phyxminiproblems-problemphyxmini0308-satisfiesexactsideonsourcegeometry, def:physics:phyx-mini-0308:phyxminiproblems-problemphyxmini0308-satisfieswaterdelaylaw, def:physics:phyx-mini-0308:phyxminiproblems-problemphyxmini0308-satisfiesaircalibrateddelaydistancelaw, def:physics:phyx-mini-0308:phyxminiproblems-problemphyxmini0308-satisfiesinterpretedbearinggeometry}
  The governing physical and auditory laws. The current numerical conclusion does not occur in any field. In particular, the far-field law retains an explicit bounded remainder
\end{definition}
\begin{proof}
  This is the declared model definition of Satisfies Sound Localization Physics.
\end{proof}
\begin{definition}[Answer Choice]
  \label{def:physics:phyx-mini-0308:phyxminiproblems-problemphyxmini0308-answerchoice}
  \lean{PhyXMiniProblems.ProblemPhyXMini0308.AnswerChoice}
  Labels of the four answer choices in the source problem
\end{definition}
\begin{proof}
  This is the declared model definition of Answer Choice.
\end{proof}
\begin{definition}[answer Choice Degrees]
  \label{def:physics:phyx-mini-0308:phyxminiproblems-problemphyxmini0308-answerchoicedegrees}
  \lean{PhyXMiniProblems.ProblemPhyXMini0308.answerChoiceDegrees}
  \uses{def:physics:phyx-mini-0308:phyxminiproblems-problemphyxmini0308-answerchoice}
  Degree readout printed beside each displayed answer choice
\end{definition}
\begin{proof}
  This is the declared model definition of answer Choice Degrees.
\end{proof}
\begin{definition}[recorded Dataset Answer]
  \label{def:physics:phyx-mini-0308:phyxminiproblems-problemphyxmini0308-recordeddatasetanswer}
  \lean{PhyXMiniProblems.ProblemPhyXMini0308.recordedDatasetAnswer}
  \uses{def:physics:phyx-mini-0308:phyxminiproblems-problemphyxmini0308-answerchoice}
  Dataset answer retained as metadata, not as a physics premise
\end{definition}
\begin{proof}
  This is the declared model definition of recorded Dataset Answer.
\end{proof}
\begin{definition}[Is Closest Angle Choice]
  \label{def:physics:phyx-mini-0308:phyxminiproblems-problemphyxmini0308-isclosestanglechoice}
  \lean{PhyXMiniProblems.ProblemPhyXMini0308.IsClosestAngleChoice}
  \uses{def:physics:phyx-mini-0308:phyxminiproblems-problemphyxmini0308-answerchoice, def:physics:phyx-mini-0308:phyxminiproblems-problemphyxmini0308-answerchoicedegrees}
  A displayed choice is uniquely closest to the interpreted degree readout
\end{definition}
\begin{proof}
  This is the declared model definition of Is Closest Angle Choice.
\end{proof}
\begin{lemma}[interpreted Bearing sine eq speed Ratio]
  \label{lem:physics:phyx-mini-0308:phyxminiproblems-problemphyxmini0308-interpretedbearing-sine-eq-speedratio}
  \lean{PhyXMiniProblems.ProblemPhyXMini0308.interpretedBearing_sine_eq_speedRatio}
  \uses{def:physics:phyx-mini-0308:phyxminiproblems-problemphyxmini0308-speedinmeterspersecond, def:physics:phyx-mini-0308:phyxminiproblems-problemphyxmini0308-soundlocalizationsetup, def:physics:phyx-mini-0308:phyxminiproblems-problemphyxmini0308-matchesintendedsideonarrival, def:physics:phyx-mini-0308:phyxminiproblems-problemphyxmini0308-hasphysicalparameters, def:physics:phyx-mini-0308:phyxminiproblems-problemphyxmini0308-satisfiessoundlocalizationphysics}
  For exact side-on source geometry, eliminating D and the common delay gives the dimensionless air-to-water speed ratio. The bounded far-field relation is not replaced by an exact plane-wave identity in this derivation
\end{lemma}
\begin{proof}
  The conclusion follows from the governing relations and figure readouts named in the statement of interpreted Bearing sine eq speed Ratio.
\end{proof}
\begin{lemma}[interpreted Bearing eq arcsin speed Ratio]
  \label{lem:physics:phyx-mini-0308:phyxminiproblems-problemphyxmini0308-interpretedbearing-eq-arcsin-speedratio}
  \lean{PhyXMiniProblems.ProblemPhyXMini0308.interpretedBearing_eq_arcsin_speedRatio}
  \uses{def:physics:phyx-mini-0308:phyxminiproblems-problemphyxmini0308-speedinmeterspersecond, def:physics:phyx-mini-0308:phyxminiproblems-problemphyxmini0308-soundlocalizationsetup, def:physics:phyx-mini-0308:phyxminiproblems-problemphyxmini0308-matchesintendedsideonarrival, def:physics:phyx-mini-0308:phyxminiproblems-problemphyxmini0308-hasphysicalparameters, def:physics:phyx-mini-0308:phyxminiproblems-problemphyxmini0308-satisfiessoundlocalizationphysics}
  The principal apparent bearing is the arcsine of the speed ratio
\end{lemma}
\begin{proof}
  The conclusion follows from the governing relations and figure readouts named in the statement of interpreted Bearing eq arcsin speed Ratio.
\end{proof}
% --- Archon declaration topology end ---
% --- Archon physics formalization source end ---
